\honest{Say Not ``Complexity''}{Eliezer Yudkowsky}{2007}

This is a story from when I first met Marcello, a young competitor in mathematics olympiads whom I was considering as an apprentice. I asked him how an AI might work out for itself how to solve a Rubik's Cube. The AI would not be preprogrammed. It would discover the cube's laws and invent the idea of an ``operator,'' a reusable sequence of moves, which is the key to solving it. He said, ``Well, I think the AI needs complexity to do X, and complexity to do Y---''

``Don't say `complexity','' I said, and explained that complexity should never be a goal in itself. \nb{Marcello had not said it was a goal.} I was thinking of people who said the Internet would ``wake up'' once it was complex enough.

He answered, ``But there's got to be some amount of complexity that does it.'' To me, saying ``complexity'' felt like the wrong move. It took a moment to find words for why, since we cannot deliberate in words over every step of our thinking. I asked whether he had read my essay on technical explanation. He had. I told him: ``Saying `complexity' doesn't concentrate your probability mass.'' He saw it at once: ``like `emergence.''' Now he had to think about how X might actually happen. ``Maybe this one is teachable,'' I thought.

\nb{The point is old and simple: a word that names a gap does not fill it. Sidney Harris drew it in a cartoon, collected in 1977, of a proof on a blackboard with the line ``then a miracle occurs'' in the middle. The post quotes that line later.}

Complexity is a useful idea in its place, I concede, with proper mathematical definitions. Even informally, you must judge whether a hypothesis is too complicated for its evidence, or how to simplify a design. The mistake is using the word to skip over something you do not understand. This mistake is ``an extremely common misstep, at least in my field,'' made in discussions of AI ``left and right, over and over again.'' I name no one making it, and say it is how human beings think by default.

It happens ``below the level of words,'' in ``a microthought.'' The cure is to notice the gaps: ``You have to feel which parts of your map are still blank, and more importantly, pay attention to that feeling.''

In a footnote I admit that the parts you can explain may turn out not to matter, though ``Mostly, people don't dare to enter terra incognita at all.''

The pressure to hide gaps, I suspect, is huge in academia. There a seemingly complete model with some ``emergent phenomena'' earns more credit than a map labelled ``then a miracle occurs,'' which a journal may not even accept. It is ``even huger'' in AI startups, where admitting that you do not know how to build AI would make your life plans collapse.

If you want examples of the skip-over, watch people discuss religion, philosophy or any science they were not trained in. Marcello and I agreed to say ``magic'' for whatever we did not understand, as a reminder that the problem was still unsolved. ``Complexity'' and ``emergence'' create an illusion of understanding. But ``magic'' leaves an honest placeholder.
